\section{Open Questions}
The following five questions remain genuinely open after this replication. They are grounded
in a re-read of the paper and this replication's execution, and each is coupled to concrete
next steps.

\begin{enumerate}
\item \textbf{Is $(K_{\text{rep}}, X)$ identifiable from $S(D)$ alone, or is there a degeneracy?}\\
\emph{Basis.} The paper fits $(K_{\text{dam}}, K_{\text{rep}}, X)$ jointly against $S(D)$ using
an independently-measured $F(D)$. But $K_{\text{rep}}$ enters $Q_1$ via $\exp(-K_{\text{rep}} P)$
and $X$ enters $Q_2$ via $P^X$ --- both act on the same monotone $P(D)$ transform. Without a
profile-likelihood or Fisher-information analysis, we cannot distinguish ``weaker repair fidelity
+ gentler direct kill'' from ``stronger repair fidelity + steeper direct kill''. Table~2
dominance assignments are a downstream reading of the fitted model, not an independent
measurement, so if the fit is degenerate the mechanism assignment is arbitrary within that
manifold.\\
\emph{Next steps.} Once Krisko \& Radman 2010 $F(D)$ and $S(D)$ are digitized: (a) re-fit with
SciPy \texttt{dual\_annealing} and record the full posterior via \texttt{emcee} or
\texttt{nautilus}; (b) plot profile likelihoods for $(K_{\text{rep}}, X)$ marginally and
jointly; (c) report the condition number of the Fisher information matrix at the MLE;
(d) if degenerate, either add a physical constraint (e.g.\ bounded $X$ per proteomic
damage-per-target reasoning) or downgrade the dominance-map claim to ``consistent with''
rather than ``demonstrates''.

\item \textbf{Does the model predict correctly on cell-death-mode-specific datasets, e.g.\ systems where SOS/RecA-driven prophage lysis dominates?}\\
\emph{Basis.} $S(D)$ is a scalar viability endpoint. The 5 rows do not stratify by death
mode. $\lambda$~IC in \textit{E.~coli} is special --- $Q_1$ is set to 1 by biological argument,
so death is entirely attributed to $Q_2$. But $\lambda$ induction is itself an SOS response,
i.e.\ downstream of DNA damage sensing --- so the ``pure $Q_2$'' assignment for $\lambda$~IC
is a definitional choice, not an empirical decomposition. The model's structural claim
(protein damage is a separate kill channel independent of DNA-damage-triggered pathways) has
never been tested against a system that isolates one from the other.\\
\emph{Next steps.} Look for RecA-independent \textit{Deinococcus} mutants (\textit{pprA$^-$},
\textit{ddrA$^-$}) and re-fit; test on \textit{Pyrococcus furiosus} or \textit{Shewanella
oneidensis} (Daly et al.\ 2007 \emph{PLoS Biol.\ }has partial data); design a synthetic experiment
where H$_2$O$_2$ or paraquat delivers pure protein oxidation with minimal DSBs and see whether
$Q_2$ alone predicts survival.

\item \textbf{How does the model behave in the low-dose regime, and is it consistent with LNT-style extrapolation?}\\
\emph{Basis.} The paper's fit spans 0--20 kGy for \textit{D.r.} and 0--4 kGy for \textit{E.c.}
--- far above bacterial doubling-inactivation regimes and well above any relevant low-dose
regime. At $D \to 0$, $P(0) \to 1$, so $Q_2 = 1$ and $Q_1 \to \exp(-K_{\text{dam}} D)$, i.e.\
pure exponential. This is a ``zero-repair-benefit'' asymptote, which is unphysical: at very
low dose, repair should be \emph{more} effective, not less. The model may thus mispredict at
low dose --- exactly the LUCID-100 regime of interest for normal-tissue late effects.\\
\emph{Next steps.} Add a low-dose repair-hyper-fidelity term
($K_{\text{dam}}^{\text{eff}} = K_{\text{dam}} (1 - r \exp(-D/D_{\text{char}}))$, $r \in [0,1]$)
and re-fit; compare to published bacterial survival curves in the 0.01--0.1~kGy range
(Battista lab \textit{D.r.} curves). If the current model over-predicts kill at low dose,
extension to eukaryotic normal-tissue late effects fails without modification.

\item \textbf{Can the framework be extended to include bystander effects, and does that reproduce mammalian shoulder-flattening?}\\
\emph{Basis.} The model is single-cell / single-hit: $D$ is per-cell, $F(D)$ is per-cell,
$Q_1$ and $Q_2$ are per-cell channels. Bystander effects (dose-independent additive
contribution from signalling between hit and unhit cells) are fundamentally incompatible with
the current multiplicative $Q_1 \cdot Q_2$ form. Any LUCID-100 extension to mammalian
systems (which show robust bystander effects at low dose, e.g.\ Little lab 3D microbeam
data) needs to answer this. The paper does not attempt it and no LUCID replication has extended
it.\\
\emph{Next steps.} Formulate $S_{\text{total}}(D) = f_{\text{hit}} (Q_1 Q_2)(D) +
f_{\text{by}} B(D, \sigma)$; fit to microbeam single-cell hit data
(Nagasawa \& Little 1999; Sawant et al.\ 2001); test whether the fitted $B$ is consistent with
a carbonyl-signalling model (bystander damage is itself protein oxidation) --- a strong
unification claim if it holds.

\item \textbf{How does the model translate to eukaryotic normal-tissue late-effects prediction?}\\
\emph{Basis.} The paper interprets $X$ as an effective count of essential protein targets whose
oxidation individually kills the cell (from the $P^X$ = product-of-independent-hits reading).
\textit{D.r.} has $X \approx 3.88$, \textit{E.c.} and $\lambda$ both $\approx 6.76$ ---
interpreted as \textit{D.r.} having a smaller vulnerable-target proteome (higher redundancy).
If this scales to eukaryotes, then tissue-specific $X$ would predict tissue-specific
late-effect susceptibility. But eukaryotic proteomes are vastly larger with more redundancy and
stronger repair --- the naive scaling would say $X$ should be much \emph{larger}, making $Q_2$
more sensitive to $P$, not less. No published test.\\
\emph{Next steps.} Refit on published mammalian survival curves for endothelial, renal, and
lung cells (Puck \& Marcus 1956 and modern colony-formation assays); check whether fitted $X$
values group by tissue radiosensitivity class (early vs late-responding); if the ordering is
inverted (more radioresistant tissues need \emph{smaller} $X$), the naive proteomic-vulnerability
interpretation fails and the DNA-and-protein framework needs re-derivation for eukaryotes.
This is the direct LUCID-100 translational question and remains fully open.
\end{enumerate}
